\section{Introduction}\label{sec:intro}

Each coefficient of a discrete Fourier transform is a weighted sum of
all the input samples. Computing the sums separately takes quadratic
work. The fast Fourier transform shares partial sums among frequencies
and, at power-of-two lengths, computes all $N$ outputs in
$O(N\log N)$ operations \cite{CooleyTukey}. How much of that
logarithmic work per sample is intrinsic? A linear circuit may choose
arbitrary scalar coefficients and reuse any intermediate result.

We prove that the optimal logarithmic exponent is $1/2$. A circuit
needs at least order $\sqrt{\log N}$ gates per sample, and for every
$\epsilon>0$ we construct circuits using
$O_\epsilon((\log N)^{1/2+\epsilon})$ gates per sample. The bounds hold for every finite abelian group whose order is
invertible in the field and whose characters take values there. They
also hold per sample when one circuit processes many independent
signals, even if its gates combine values from different signals.

Write $C_K(T)$ for the smallest number of gates in an acyclic linear
circuit computing a matrix $T$ over a field $K$. Each gate forms
$au+bv$ from two available values, with arbitrary $a,b\in K$. An
output may select and scale a single input or gate value, or return
zero. All logarithms in the paper are to base two.

\begin{theorem}[The Fourier exponent]\label{thm:main}
Let $G_1,\ldots,G_J$ be finite abelian groups of common order $N\ge2$,
where $J\ge1$. Suppose $N$ is invertible in $K$ and $K$ contains a
primitive root of unity of order $\exp(G_j)$ for each $j$. Let $T$ be
the direct sum of their Fourier transforms, choosing the forward or
inverse transform independently in each block. For an absolute $c>0$
and a constant $C_\epsilon$ depending only on $\epsilon>0$,
\[
 c\,JN\sqrt{\log N}\ \le\ C_K(T)\
            \le\ C_\epsilon\,JN(1+\log N)^{1/2+\epsilon}.
\]
\end{theorem}

\needspace{4\baselineskip}
Here $\exp(G)$ is the exponent of the group. The field condition gives
one frequency for each character $\chi:G\to K^\times$. The forward
transform sends a list $(x_g)_{g\in G}$ to the list of sums
$\sum_{g\in G}\chi(g)x_g$, one for each $\chi$. For example, on
$(\mathbb Z/2\mathbb Z)^k$ over $\mathbb C$ these are the
Walsh--Hadamard sums, with coefficients $+1$ and $-1$. Each of $J$
independent Walsh--Hadamard transforms of length $2^k$ costs at least
$c\,2^k\sqrt{k}$ gates on average, however fast $J$ grows with $k$.
Theorem~\ref{thm:main} leaves open an upper bound with $\epsilon=0$.

The exponent also governs \emph{formal polynomial convolution}, the
map that takes two lists of $n$ coefficients to the coefficients of
their polynomial product. Over every field, $J$ independent products
require $\Omega(Jn\sqrt{\log n})$ gates in a division-free arithmetic
circuit, and $O_\epsilon(Jn(1+\log n)^{1/2+\epsilon})$ gates suffice
(Corollary~\ref{cor:conv-exponent}). Here addition, subtraction and
multiplication are binary operations, including multiplication by a
prescribed scalar. For the Fourier transform itself, allowing products
of intermediate polynomials in an exact formal circuit also gives the
same exponent. Taking the linear part of each gate turns a product
$uv$ into $u_0v_1+v_0u_1$, where the subscripts denote constant and
linear parts. Conversely, three field operations suffice for a weighted
linear gate.

\subsection{A lower bound with feedback and errors}

The lower bound allows two further ways in which a circuit might save
work. First, it may use feedback. A system of $s$ scalar states
specifies each state as a weighted sum of at most two inputs or
states, with a unique solution for every input. Write $\mathcal C_s(h)$
for the $h$-input, $h$-output matrices computed in this way over
$\overline K$ using at most $s$ states and the readout convention
above. Second, the system may leave behind a low-rank error and an
error concentrated on a prescribed number of matrix entries. Each
can couple all the blocks. Even a limit of these corrected systems
obeys the lower bound.

\begin{theorem}[Robust lower bound]\label{bch:uniform-robust}
With $T$ as in Theorem~\ref{thm:main} and $h=JN$, there are absolute
constants $c,\beta,\eta>0$ with the following property. Suppose $T$
belongs to the Zariski closure over $\overline K$ of all sums $M+B+Z$
satisfying
\[
 M\in\mathcal C_s(h),\qquad \operatorname{rank}B\le\beta h,\qquad
 |\operatorname{supp} Z|\le\eta h\,2^{\sqrt{\log N}/8}.
\]
Then $s\ge c\,h\sqrt{\log N}$.
\end{theorem}

The support budget counts nonzero entries in the entire $h\times h$
matrix. Taking the closure of the sums lets the feedback matrix and
both corrections vary together. Setting the corrections to zero and
using an acyclic circuit gives the lower bound of Theorem~\ref{thm:main}.

\subsection{Earlier bounds for Fourier circuits}

The classical $O(N\log N)$ count has two ingredients. Cooley and
Tukey factor a transform of length $ab$ into transforms along the two
axes, with a diagonal scaling between them. Good's factorization
eliminates the scaling for coprime lengths
\cite{Good,CooleyTukey,cooleyLewisWelch1967}. Recursing on powers of
two gives the familiar FFT. Bluestein's reduction to convolution
with a fixed quadratic phase, followed by padding to a power of two,
gives the same asymptotic count for every complex transform length
\cite{Bluestein,rabinerSchaferRader1969}.

Several classical lower bounds explain when this count is necessary.
Papadimitriou proved optimality under restrictions on the paths in
the computation graph \cite{Papadimitriou}. Winograd counted multiplications
separately from additions \cite{Winograd1979DFT}. For the unnormalized
complex transform, the determinant has modulus $N^{N/2}$.
Morgenstern showed that if scalar multipliers have modulus at most one,
an addition can at most double the largest determinant modulus
attainable from the available forms. At least $\tfrac12N\log_2 N$
additions are therefore necessary \cite{Morgenstern}. Normalizing the
transform makes its determinant have modulus one. For this transform,
Ailon proved an $\Omega(N\log N)$ lower bound for in-place
two-coordinate unitary operations, and an $\Omega(N\log N/R)$ bound
for in-place rotations and scalings when every intermediate map has
condition number at most $R$ \cite{Ailon,AilonCondition}.
Theorem~\ref{thm:main} permits arbitrary scalar coefficients,
intermediate storage and reuse.

For power-of-two lengths, Alman and Rao used matrix nonrigidity to
reduce the leading operation counts to $\tfrac{23}{24}N\log_2 N$
for Walsh--Hadamard transforms over a field and $\tfrac{15}{4}N\log_2 N$
in real arithmetic for complex cyclic transforms, up to $O(N)$ terms
\cite{AlmanRao}. Finite tensor constructions subsequently gave
$o(N\log N)$ circuits along a subsequence \cite{FiniteFourier}.
At every length, the algorithm of \cite[Theorem~1.1]{ExplicitFourier}
uses a finite network from \cite[Proposition~8 and Section~3.5]{Tape}
to save a fixed positive power of $\log N$. It charges for array indexing and for preparing the
scalars from a supplied root of unity of polynomial order.
Sussman developed a recursion that handles several target maps
together \cite{Sussman}. Kartal combined it with constructions of
CubeJerry, Cinna Davis and Ali Romman, and Yao and Liu refined the
resulting numerical bounds \cite{Kartal,kartalpi,YaoLiu}. A subsequent
construction gives a saving greater than $0.0039$ in this complex
arithmetic model \cite{RelayFourier}. The circuit size in Theorem~\ref{thm:main}
charges only for the binary linear gates.

Valiant related linear circuits to matrix rigidity, the number of
entries that must change to reach a prescribed rank
\cite{Valiant1975,Valiant1977}. Alman and Williams showed that the
Walsh matrix fails Valiant's rigidity criterion \cite{AlmanWilliamsRigidity}.
For every fixed $\delta>0$, Dvir and Liu reduce the rank of any
sufficiently large complex abelian character table to
$o(N/\log\log N)$ by changing at most $N^{1+\delta}$ entries
\cite[Theorem~6.1]{DvirLiu}. Theorem~\ref{bch:uniform-robust},
with one transform and no circuit states, gives a lower bound in the
same variables. Reducing the rank to $\beta N$ requires more than
$\eta N2^{\sqrt{\log N}/8}$ changes. For a smaller target rank, Alman
showed that reducing a Walsh matrix over $\mathbb F_3$ to rank
$\lfloor(\log_2 N)^2/80\rfloor$ requires changing more than one percent
of its entries \cite{AlmanRigidity}.

For convolution, specializing one input turns a bilinear circuit into
a linear one. Raz used specialization for bounded-coefficient matrix
multiplication \cite{Raz}. B\"urgisser and Lotz proved an
$\Omega(n\log n)$ bound for complex cyclic convolution and ordinary
polynomial multiplication in the bounded-coefficient bilinear model
\cite{BurgisserLotz}.
Corollary~\ref{cor:conv-exponent} uses a fixed operand over the algebraic
closure to obtain a Fourier transform, giving a lower bound for
division-free formal polynomial multiplication with arbitrary scalars
over every field.


\subsection{Where the square root enters}

\paragraph{A geometric test for a small circuit.}
Rank and the number of nonzero entries each describe one way a matrix
can be easy to compute. Neither by itself explains the Fourier bound.
For example, over $\mathbb Q$ the map that adds the sum of all inputs
to each input is dense and invertible, yet has a linear-size circuit.
Part~\ref{part:lower} tests rank and support together by placing the
indices of a matrix at the corners of a rectangle. The construction
uses the tensor prices of the companion geometry paper \cite{oaG}.
A wire sharing one of $q$ labels between two sites costs their
distance times $\log q$. Placing the row label at both left corners
and the column label at both right corners gives a price that reads
log rank when the rectangle has height zero and log support size when
it has width zero, per unit of the remaining length.

At width one and height $L$ the price of a character table is exactly
$(L+\sqrt{1+L^2})\log N$ (Theorem~\ref{geo:saturation}). A
permutation matrix at the same placement costs $(L+1)\log N$. Both
matrices have full rank. As the rectangle opens up, their prices
separate by approximately $\tfrac12 L^2\log N$. Thus even at height
of order $1/\sqrt{\log N}$ the test detects a difference of constant
size. To obtain the exact Fourier price, we use two diagonal
character phases and changes of basis to exchange the two labels.
The cost of this exchange forces the character table to attain the
upper bound from supplying its labels through the center.

The circuit comparison applies the rectangle test to the list of both
inputs and outputs of a transform. Pairing two copies of the list
exposes the character table, and the resulting lower bound on the
list survives a small fractional loss of rank. For a batch, eliminating
a single rank correction that mixes the signals might seem to destroy
every individual list. Row reduction and a scaling limit recover a
separate projection on each block. Their rank losses add up to at
most the original rank budget, so most blocks retain the lower bound.

Finally, repeat each scalar input and state along $d$ time coordinates.
A shift at every state leaves a copy of the desired transfer on each
time diagonal. Changing the shift by a map of rank one kills every
path that runs through $d$ states. The short paths left behind reach
at most $2^{d-1}$ inputs per output, while the total rank correction
costs at most one dimension per state. With $d$ of order
$\sqrt{\log N}$, the list test rules out the resulting combination
of rank and support unless the circuit has at least a constant times $JN\sqrt{\log N}$
states (Section~\ref{bch:proof:section}).

\paragraph{Predicting a product and repairing it exactly.}
The upper bound in Part~\ref{part:finite} starts with the Walsh--Hadamard
transform. We express a transform on many binary digits, up to
permutations and scalings, as a product of $m$ commuting operators.
In their common frequency basis each operator multiplies by a sign.
The product is determined by the parity of the number of negative
signs. Run many independent inputs together and assign them all
possible sign shifts in equal proportions. At each frequency the
shifted count has a binomial distribution over the batch. Almost all
counts lie in an interval of width $O(\sqrt{m\log m})$, on which a
polynomial of that degree gives the parity sign exactly. Undoing the
known shift gives the desired answer on those rows.

The exceptional rows depend on the frequency. Converting every input
to the common basis to correct them would spend the work we are trying
to save. Instead, send a small number of linear combinations of the
rows into the frequency basis. There the exceptional positions are
known, and interpolation recovers the corrections from those combinations.
A fixed number of further predictor applications in ordinary
coordinates then corrects every output (Theorem~\ref{thm:repair}).

A degree bound for the predictor still has to become a gate bound.
We evaluate it in $O(\sqrt{m\log m})$ stages, each a matrix linear
in the commuting operators. A stage may ask for short Fourier
transforms in $m$ different directions. The construction in
Section~\ref{sec:prime-kernel} shares those calls among rows by
factoring their routes through common source and destination maps.
It uses different characteristics for the routing labels and the
addresses, drawing on constructions in \cite{Tape,FranklWilson1981,Alon1998}
and related ideas in index coding \cite{BarYossefBirkJayramKol2006,LubetzkyStav2009}.
For any fixed $\theta>1/2$, a finite choice of parameters then
replaces a transform on $Af$ digits by fewer than $A^\theta$
transforms on $f$-digit axes per input row, averaged over the batch.
A single large transform already contains the independent inputs this
step needs. Holding some digits fixed gives disjoint slices on which
the remaining digits are transformed, and these slices fill the batches
in the recursion. Iterating gives the exponent on prime tensor powers.
In the binary
case the construction also controls how errors in individual gates
affect the normalized transform (Proposition~\ref{prop:binary-template}).
The tape multiplication paper uses this implementation \cite{oaT}.

\paragraph{Returning to the given field.}
A general finite abelian group factors into cyclic groups of different
orders. Even when the field contains every character of the original
group, a fast algorithm may call for roots of unity at other lengths.
Part~\ref{part:all} packs independent data into the coordinates of
small field extensions and charges the operations there back to the
base field. For formal convolution in characteristic zero, a single
deformation parameter handles all cyclic factors at once. In positive
characteristic, finite-field extensions supply roots with short
Frobenius periods. The resulting multiplication circuits over the
original field give fast evaluation and interpolation on geometric
grids, of which every cyclic Fourier matrix is an instance
(Theorems~\ref{thm:conv} and~\ref{up:grid}).
In the reverse direction, fixing one operand of a polynomial product
realizes a cyclic Fourier transform after diagonal scalings. Applying
the Fourier lower bound to this specialization gives the convolution
lower bound in Corollary~\ref{cor:conv-exponent}.
